% chapters/10/praktikum.tex
\begin{praktikumbox}
Praktikum ini mengikuti Tugas Praktikum 08 pada RPS: memvisualisasikan A* pada
peta berpetak, dengan titik awal, titik tujuan, dan dinding yang ditandai,
serta jalur yang dianimasikan.

\begin{langkah}
  \item \textbf{Jalankan programnya.} Ketik ulang \code{ch10_astar.py} dari
        Daftar~\ref{lst:ch10-astar} ke dalam folder \code{code}. Jalankan, lalu
        amati penyebaran pencarian dari titik awal menuju titik tujuan.

  \item \textbf{Bandingkan ketiga heuristik.} Jalankan program dengan pilihan
        \code{--bandingkan}, lalu catat tabel yang dihasilkan.
\begin{kodekeluaran}
python ch10_astar.py --bandingkan
\end{kodekeluaran}

  \item \textbf{Amati perbedaan penyebarannya.} Jalankan program dengan
        \code{--heuristic manhattan}, lalu tekan tombol H untuk berganti
        heuristik. Amati bagaimana bentuk penyebaran pencarian berubah, dan
        bagaimana jumlah petak yang diperiksa bertambah.

  \item \textbf{Ubah petanya.} Pindahkan huruf \code{G} ke tempat lain, lalu
        tambahkan beberapa dinding. Perhatikan bagaimana jalurnya berubah.
        Buatlah satu keadaan yang membuat tujuan tidak dapat dicapai, lalu
        amati keluaran program pada terminal.

  \item \textbf{Simpan tangkapan layar.} Jalankan program dengan pilihan
        \code{--simpan-gambar} untuk menyimpan keadaan pencarian sebagai
        berkas gambar, lalu sertakan gambar itu pada laporan praktikum Anda.
\begin{kodekeluaran}
python ch10_astar.py --heuristic manhattan --simpan-gambar hasil.png
\end{kodekeluaran}

  \item \textbf{Uji heuristik yang melebihkan jarak.} Gandakan hasil
        \code{jarak_manhattan} pada fungsi heuristiknya. Jalankan
        \code{--bandingkan}, lalu bandingkan biaya jalur dengan hasil
        sebelumnya. Jelaskan apa yang terjadi.

  \item \textbf{Tulis catatan praktikum} yang memuat tabel perbandingan
        heuristik, tangkapan layar pencarian, dan analisis pengaruh heuristik
        terhadap kinerja dan terhadap mutu jalur.

  \item \textbf{Simpan pekerjaan} ke repositori Git dengan pesan penyimpanan
        yang menerangkan isinya.
\end{langkah}
\end{praktikumbox}

\begin{table}[htbp]
\centering
\caption{Lembar kerja Praktikum 08}
\label{tab:lembar-praktikum-08}
{\setlength{\tabcolsep}{4pt}\footnotesize%
\begin{tabularx}{\textwidth}{@{}C{3.0cm} C{2.6cm} C{2.4cm} Y@{}}
\toprule
\headrow \textbf{Heuristik} & \textbf{Petak diperiksa} & \textbf{Biaya jalur} & \textbf{Bentuk penyebaran pencarian} \\
\midrule
Manhattan & & & \\[1.4em]
\addlinespace
Euclidean & & & \\[1.4em]
\addlinespace
Tanpa perkiraan & & & \\[1.4em]
\addlinespace
\multicolumn{4}{@{}p{\textwidth}@{}}{\footnotesize\itshape
Hasil pada heuristik yang digandakan (langkah keenam): biaya jalur \dotfill} \\[1.4em]
\addlinespace
\multicolumn{4}{@{}p{\textwidth}@{}}{\footnotesize\itshape
Keadaan ketika tujuan tidak dapat dicapai: \dotfill} \\[1.4em]
\bottomrule
\end{tabularx}}
\end{table}

\begin{catatabox}
Langkah keenam adalah langkah yang paling penting bagi pemahaman indikator
\ind{b}. Heuristik yang digandakan akan membuat pencarian tampak lebih cepat,
padahal jalur yang dihasilkannya dapat lebih panjang. Program
tidak akan melaporkan kesalahan apa pun; hanya angka biaya jalur yang
menunjukkan adanya masalah.
\end{catatabox}
